\section{Predictive phase continuation}

Let $\mathcal O_{\rm cal}$ denote registered calibration observables and let
$\mathcal O_{\rm solv}(\bm\theta)$ be a solvable surrogate.  We estimate
\begin{equation}
 \widehat{\bm\theta}_{\rm eff}=\arg\min_{\bm\theta}
 D\!\left(\mathcal O_{\rm cal}^{\rm target},
 \mathcal O_{\rm solv}(\bm\theta)\right),
\end{equation}
then predict $\gchat$, the order curve, and the response to an intervention.
The bridge error is evaluated only on withheld quantities,
\begin{equation}
 \Ebridge=D\!\left(\mathcal O_{\rm holdout}^{\rm target},
 \mathcal O_{\rm solv}(\widehat{\bm\theta}_{\rm eff})\right).
\end{equation}
For each edge we also record semantic retention
$S_{\rm sem}=I(m;Z)/H(Z)$ and normalized predictive error
$|\gc-\gchat|/(\sigma_{\gc}+\epsilon)$.

Two assumptions need not combine additively.  Their interaction is
\begin{equation}
 \Delta_{ij}\gc=\gc^{ij}-\gc^i-\gc^j+\gc^0.
\end{equation}
The operational outcome vocabulary includes stable and renormalized
continuation, splitting, merging, rounding, new-phase insertion,
computational--statistical separation, path dependence, and semantic failure.
These labels are not inferred from visual smoothness.

We distinguish macrostate perplexity $K_{\rm eff}^{\rm macro}=\exp H(Z)$ from
correlation-based independent components and spectral effective rank.  Each
PhaseCard stores estimator, probe distribution, normalization, units,
finite-size variable, intervention, and null model.
